% TOP_PROBLEM: {"id": 2302076, "title": "Research Problems in Function Theory — Problem 2.76", "category_id": 8, "category": {"id": 8, "name": "analysis", "display_name": "Analysis", "description": "Limits, continuity, calculus, and function theory.", "slug": "analysis", "order_index": 8, "created_at": "2026-07-31T15:26:25.671Z"}, "status": "open"}
% FIRST_POSTED: null
% ENTRY_KIND: statement_only
% Imported from: batch12.tex; UnsolvedMath ID 2302076
% Compile twice: lualatex 2302076.tex
\documentclass[11pt]{article}
\usepackage{fontspec}
\setmainfont{Latin Modern Roman}
\usepackage{amsmath,amssymb,mathtools,unicode-math}
\setmathfont{Latin Modern Math}
\usepackage{xurl,hyperref}
\hypersetup{colorlinks=true,urlcolor=blue,pdftitle={UnsolvedMath Problem 2302076}}
\newcommand{\sref}[2]{\hyperlink{source:#1:#2}{[S#2]}}
\newcommand{\cataloguescope}[1]{\paragraph{Mathematical question.}#1}
\begin{document}
% UnsolvedMath ID 2302076; source catalogue code AMR-022-2076
\setcounter{section}{2302075}
\section{Research Problems in Function Theory — Problem 2.76}\label{2302076}
{\small\sffamily UnsolvedMath ID 2302076 · Analysis}
\subsection{Definitions and mathematical statement}

\paragraph{Shared notation.}$\mathbb N=\{1,2,\ldots\}$, and $\mathbb Z,\mathbb Q,\mathbb R,\mathbb C$ denote the integers, rationals, reals and complex numbers. Any other conventions are specified locally.


For an entire function $f$, write $f^{\circ0}(z)=z$ and $f^{\circ(n+1)}=f\circ f^{\circ n}$. Its Fatou set is the set on which the iterates form a normal family: every sequence has a subsequence converging locally uniformly in the spherical metric to a meromorphic function or to $\infty$. A Fatou component is a connected component of this set. For a set $E\subset\mathbb C$, $\dim_H E$ denotes its Hausdorff dimension, defined by the critical exponent of the Hausdorff measures $\mathcal H^s(E)$.
The source update describes counterexamples of Bishop: boundaries of dimension $1$ need not be circles or lines.
\paragraph{Statement recorded in the supplied source.}
For every entire function $f$ and every Fatou component $\Omega$, is it true that
\[\dim_H\partial\Omega>1\quad\text{or}\quad\partial\Omega\text{ is a Euclidean circle or a straight line}?\] \sref{2302076}{2}
\subsection{Short English statement}
Must the boundary of an entire-function Fatou component have Hausdorff dimension greater than one unless it is a circle or a line?
\subsection{Sources}
\begin{description}
\item[\hypertarget{source:2302076:1}{S1}] ULAM.AI and the UnsolvedMath Contributors, UnsolvedMath: A Curated Collection of Open Mathematics Problems, record 2302076 (AMR-022-2076). CC BY 4.0. \url{https://huggingface.co/datasets/ulamai/UnsolvedMath}.
\item[\hypertarget{source:2302076:2}{S2}] W. K. Hayman and E. F. Lingham, \emph{Research Problems in Function Theory}, arXiv:1809.07200v2 (2018), Problem 2.76 and Update 2.76. \url{https://arxiv.org/abs/1809.07200}.
\end{description}
\label{end:2302076}
\end{document}
